\documentclass[11pt]{article}

\usepackage[T1]{fontenc}
\usepackage[utf8]{inputenc}
\usepackage{lmodern}
\usepackage[a4paper,margin=28mm]{geometry}
\usepackage{microtype}
\usepackage{amsmath,amssymb,amsthm,mathtools}
\usepackage{booktabs}
\usepackage{tabularx}
\usepackage{longtable}
\usepackage{array}
\usepackage{enumitem}
\usepackage{xcolor}
\usepackage{listings}
\usepackage{xurl}
\usepackage{hyperref}

\definecolor{linkblue}{RGB}{30,70,130}
\definecolor{codebg}{RGB}{247,248,250}
\definecolor{codeframe}{RGB}{210,214,220}
\hypersetup{
  colorlinks=true,
  linkcolor=linkblue,
  citecolor=linkblue,
  urlcolor=linkblue,
  pdftitle={Higher Abductive Theory Discovery},
  pdfauthor={Concept paper}
}

\lstdefinestyle{spec}{
  basicstyle=\ttfamily\footnotesize,
  backgroundcolor=\color{codebg},
  frame=single,
  rulecolor=\color{codeframe},
  breaklines=true,
  columns=fullflexible,
  keepspaces=true,
  showstringspaces=false,
  xleftmargin=0.5em,
  xrightmargin=0.5em,
  aboveskip=0.7em,
  belowskip=0.7em
}

\setlist[itemize]{leftmargin=1.6em,itemsep=0.25em,topsep=0.35em}
\setlist[enumerate]{leftmargin=1.8em,itemsep=0.3em,topsep=0.35em}
\renewcommand{\arraystretch}{1.18}

\newtheorem{definition}{Definition}[section]
\newtheorem{principle}[definition]{Principle}
\newtheorem{proposition}[definition]{Proposition}
\newtheorem{remark}[definition]{Remark}

\newcommand{\Th}{\mathbf{Th}}
\newcommand{\Cat}{\mathbf{Cat}}
\newcommand{\Mod}{\mathsf{Mod}}
\newcommand{\Obs}{\mathsf{Obs}}
\newcommand{\Scope}{\mathsf{Scope}}
\newcommand{\Evidence}{\mathsf{Evidence}}
\newcommand{\Rel}{\mathsf{Rel}}
\newcommand{\Press}{\mathsf{Press}}
\newcommand{\Comp}{\mathsf{Comp}}
\newcommand{\Gen}{\mathsf{Gen}}
\newcommand{\Val}{\mathsf{Val}}
\newcommand{\Cl}{\mathsf{Cl}}
\newcommand{\Def}{\mathsf{Def}}
\newcommand{\Cost}{\mathsf{Cost}}
\newcommand{\id}{\mathrm{id}}

\title{\textbf{Higher Abductive Theory Discovery}\\
\large A Categorical and Homotopical Framework for Relevance-Guided Generative Reasoning}
\author{Conceptual research programme}
\date{July 2026}

\begin{document}
\maketitle

\begin{abstract}
This paper develops a framework for systematic theory discovery across mathematics and physics.
The central proposal is to represent partial theories, models, observations, and transformations as
a higher theory atlas. Category theory supplies the grammar of composition, interpretation, and
universal construction; homotopy type theory supplies a proof-relevant treatment of identity,
equivalence, deformation, and higher coherence. Abduction is then formulated as the constrained
completion of defective diagrams: missing representing objects, noncommuting translations, failed
descent, absent adjoints, broken factorisations, and unresolved obstructions become mathematically
structured locations at which a new concept may enter.

Completion alone, however, cannot distinguish a scientifically important discovery from an elegant
but sterile mathematical construction. We therefore introduce a relevance layer. A research context
induces a field of conceptual pressures over the atlas: empirical anomalies, proof bottlenecks,
explanatory fragmentation, duplicated structure, representational instability, and pragmatic needs.
A philosophical-reflective generative layer reconstructs the meaning, inferential role, operational
anchors, and modal commitments of concepts at essential points of fragmentation before those
pressures are compiled into formal defects. A generative element then proposes typed theory
transformations. Candidates are evaluated by separate operators for unblocking relevance, explanatory
power, unification power, predictive yield, theorem transport, pragmatic value, conceptual necessity,
robustness, and concrete return. Relevance guides search but does not establish truth; proof, model
construction, simulation, counterexample search, and experiment remain independent validation
channels. The paper concludes with an implementation-grade agentic architecture: typed skills, data
contracts, workflows, validation boundaries, storage and prover adapters, evaluation procedures, and
a staged construction programme. The strongest attainable guarantee remains relative: within an
explicit theory atlas and operator grammar, the system can identify, compare, instantiate, and test
minimal high-leverage transformations, while diagnosing when progress requires a revision of the
representation space itself.
\end{abstract}

\tableofcontents
\newpage

\section{Motivation}

The discovery of a new theory is rarely a matter of unconstrained invention. It usually takes place
inside a partially mapped conceptual field. Several descriptions work in restricted regimes; recurring
calculations suggest a common structure; local models fail to assemble globally; two mature theories
make incompatible identifications; or a proof programme repeatedly reaches the same obstruction.
Creative work consists partly in recognizing what kind of missing concept would reorganize these
pressures into a more coherent whole.

There are therefore two complementary problems:
\begin{enumerate}
  \item How can a machine represent and reason over a space of theories, models, transformations,
  and equivalences without reducing it to an untyped semantic graph?
  \item How can it determine which formally available completion matters for the actual mathematical
  or physical problem?
\end{enumerate}

Category theory is well suited to the first problem because it treats relations, translations,
composition, and universal properties as primary. Homotopy type theory (HoTT) enriches this picture
by making identity proof-relevant and by retaining transformations between transformations
\cite{hottbook}. Neither discipline, by itself, answers the second problem. A category contains many
possible completions, most of which have no scientific importance. Relevance must therefore be made
an explicit component of the discovery architecture.

The thesis of this paper is:

\begin{quote}
\emph{Theory discovery can be organized as relevance-guided search for semantics-preserving
completions of structured defects in a higher theory atlas, together with controlled revision of the
atlas when its existing representational grammar becomes the obstacle.}
\end{quote}

This formulation is deliberately weaker than the claim that all conceptual discovery can be
mechanized. It does not presuppose an absolute category of every possible theory or an exhaustive
metric of explanatory value. It instead makes the scope, transformations, pressures, and evaluative
commitments explicit, so that both the power and the limits of the reasoner can be examined.

\section{The higher theory atlas}

\subsection{Theories as structured, scoped objects}

A theory used for discovery must contain more than a signature and a set of axioms.

\begin{definition}[Theory fragment]
A theory fragment is a structured object
\begin{equation}
T=(\Sigma,A,\Mod,\Obs,\Scope,\Evidence),
\end{equation}
where:
\begin{itemize}
  \item $\Sigma$ contains types, concepts, symbols, and admissible constructions;
  \item $A$ contains axioms, laws, constraints, dynamics, and inference rules;
  \item $\Mod(T)$ is the category or higher groupoid of models and model maps;
  \item $\Obs(T)$ specifies observables, measurements, computational outputs, or target judgments;
  \item $\Scope(T)$ specifies regimes, scales, idealizations, background conditions, and approximation
  status;
  \item $\Evidence(T)$ links claims to proofs, examples, counterexamples, simulations, and experiments.
\end{itemize}
\end{definition}

For mathematics, $\Obs(T)$ may contain canonical invariants, decision problems, and conjectures.
For physics, it includes the operational bridge from formal states to quantities that can be measured
or computed. Scope is essential in both cases: a low-energy effective theory, a finite-dimensional
theorem, and an asymptotic approximation must not be treated as unrestricted claims.

\subsection{Morphisms and higher cells}

Let $\Th$ denote a category, bicategory, double category, or $(\infty,1)$-category of theory
fragments. A theory morphism
\begin{equation}
f:T\longrightarrow T'
\end{equation}
interprets the vocabulary and validated structure of $T$ inside $T'$. It induces a contravariant
operation on models,
\begin{equation}
f^*:\Mod(T')\longrightarrow \Mod(T).
\end{equation}

This separates the translation of theoretical language from the migration of models. The MMT
framework already treats formal knowledge as modular theories related by theory morphisms and
mediates across different logical foundations \cite{rabe2016}. Generalized algebraic theories provide
a restricted computational setting in which theory morphisms and model migration can be implemented
declaratively \cite{gatlab}.

An ordinary directed graph is insufficient because several scientifically different relations must be
distinguished. A practical atlas can be organized as a double category:
\begin{itemize}
  \item horizontal morphisms represent translations, dualities, changes of description, and exact
  representation maps;
  \item vertical morphisms represent specialization, approximation, refinement, coarse-graining, and
  limiting regimes;
  \item squares state that a translation is compatible with a refinement;
  \item higher cells compare alternative compatibility proofs.
\end{itemize}

The atlas must also retain partiality. Two theories may overlap on only a sublanguage, a morphism may
be defined only in one regime, and a comparison may commute only up to a controlled approximation.
Forcing every relationship to be total would manufacture false unifications.

\subsection{The role of homotopy type theory}

HoTT contributes a disciplined treatment of equivalence and coherence:
\begin{enumerate}
  \item Identity types retain the paths by which two structures are identified.
  \item Higher identity types compare identifications themselves.
  \item Univalence permits equivalent structures to be identified at the appropriate level, reducing
  duplication by presentation.
  \item Higher inductive types permit free construction by generators, relations, and explicit gluing.
  \item Truncations control how much higher information is retained.
  \item Modalities can distinguish exact from approximate, local from global, classical from quantum,
  or Euclidean from Lorentzian contexts.
\end{enumerate}

HoTT should not be imposed uniformly. Many useful theory graphs can be implemented using dependent
type theory and ordinary categorical structures. Higher structure is justified when multiple paths of
translation, deformation, or gluing are themselves part of the scientific content.

\section{Abduction as defect-driven completion}

Let
\begin{equation}
D:J\longrightarrow\Th
\end{equation}
be a diagram representing a mapped collection of theory fragments. An opportunity for discovery is
not simply an empty location in a graph. It is a typed defect: a failure of a desired structural or
semantic property.

\begin{definition}[Completion problem]
A completion problem is a tuple
\begin{equation}
\mathcal{O}=(D,\delta,\mathcal{K},C),
\end{equation}
where $D$ is a theory diagram, $\delta$ is a detected defect, $\mathcal{K}$ is an admissible class of
extensions, and $C$ contains semantic, evidential, and scope constraints. A candidate solution is an
extension $D\hookrightarrow\widehat D$ in $\mathcal{K}$ that removes or explains $\delta$ while
respecting $C$.
\end{definition}

The main defect--completion operators are summarized in Table~\ref{tab:defects}.

\begin{table}[ht]
\centering
\small
\begin{tabularx}{\textwidth}{>{\raggedright\arraybackslash}p{0.20\textwidth}
>{\raggedright\arraybackslash}p{0.23\textwidth}
>{\raggedright\arraybackslash}p{0.24\textwidth}
X}
\toprule
Detected defect & Categorical question & Abductive proposal & Validation demand \\
\midrule
Recurrent natural family without a corresponding object & Is the functor representable? & New object, state, field, invariant, or space & Prove the universal property and construct models \\
Two equivalent paths disagree & Is a natural transformation, homotopy, or correction term missing? & Coupling, interaction, contextual dependency, or coherence law & Prove coherence and identify observable consequences \\
Local models agree but fail to assemble globally & What is the descent or cohomological obstruction? & Topological sector, anomaly, boundary datum, or global degree of freedom & Compute the obstruction and recover local behavior \\
Theories share an interface but remain separate & Does a pushout or homotopy pushout give a joint extension? & Unifying theory or common generative core & Show model existence and conservative recovery of inputs \\
Theorem transport fails across scale or representation & Is an adjoint, Kan extension, lift, or error object missing? & Effective variable, renormalization law, or controlled approximation & Establish transport bounds and limiting cases \\
Many derivations factor through an unnamed structure & Is there a universal factorization? & Mediating concept or organizing principle & Demonstrate uniqueness, compression, and new consequences \\
An expected exactness or lifting property fails & Which obstruction measures the failure? & New invariant or no-go principle & Prove necessity and exhibit discriminating examples \\
\bottomrule
\end{tabularx}
\caption{Defect-driven generators of candidate concepts.}
\label{tab:defects}
\end{table}

These operators make creativity constrained rather than mechanical. A pushout of signatures, for
example, is not automatically a unification: the combined theory may have no models or may destroy
the validated meaning of its components. A homotopy pushout records richer identifications, but it
still requires a semantic interpretation.

\subsection{Abductive admissibility}

Suppose $T$ is a base theory and $E$ an unexplained observation, partial model, or proof demand. An
abductive proposal is an extension
\begin{equation}
i:T\longrightarrow T_H.
\end{equation}
It is initially admissible only when:
\begin{enumerate}
  \item $\Mod(T_H)$ is nonempty;
  \item there exists $M_H\in\Mod(T_H)$ whose reduct $i^*M_H$ agrees with the validated part of $E$;
  \item $i$ is conservative over protected content, or every revision is made explicit;
  \item $T_H$ entails consequences not already encoded in $E$;
  \item at least one path exists from these consequences to concrete proofs, models, calculations,
  interventions, or experiments.
\end{enumerate}

Abduction proposes what would make the diagram work. Deduction determines what accepting the
proposal commits us to. Model construction establishes realizability. Falsification and experiment
test whether the commitment is borne by the intended world.

\section{Why completion mechanics are insufficient}

A category may contain indefinitely many completions, universal constructions, adjoints, and formal
factorizations. Most are irrelevant to the research problem. Three distinctions are therefore
fundamental.

\subsection{Truth, relevance, and novelty}

Truth and relevance are not interchangeable. A hypothesis can be highly relevant and false; a
theorem can be true and irrelevant to the target question. Novelty is also independent: a new formal
construction may be correct but explanatorily empty.

The architecture must enforce a separation of powers:
\begin{equation}
\text{generation and relevance routing}
\quad\neq\quad
\text{formal or empirical validation}.
\end{equation}

The generative component is allowed to be speculative. It is not allowed to certify its own output.

\subsection{Relevance is conditional}

There is no context-free ranking of theoretical importance. Relevance is conditional on a research
context
\begin{equation}
Q=(q,\mathcal{E},\mathcal{S},\mathcal{B},\mathcal{V}),
\end{equation}
where $q$ is the target question, $\mathcal{E}$ the accepted evidential standards, $\mathcal{S}$ the
protected scope, $\mathcal{B}$ the computational or experimental budget, and $\mathcal{V}$ the
research values and intended uses. Thus
\begin{equation}
\Rel(H\mid Q,D,E)
\end{equation}
is meaningful, whereas an absolute $\Rel(H)$ generally is not.

\subsection{Universality is relative to representation}

Every universal construction is universal relative to selected objects and morphisms. A poor atlas can
make an unimportant object appear universal or prevent an important concept from being expressible.
The reasoner must consequently distinguish completion within an atlas from revision of the atlas.

\section{Conceptual pressures}

\subsection{Pressures as typed demands}

A conceptual pressure is evidence that the present theory organization is inadequate relative to $Q$.
It is not necessarily a contradiction. It may be inefficiency, fragmentation, fragility, or unrealized
opportunity.

\begin{definition}[Pressure item]
A pressure item is
\begin{equation}
p=(s,\tau,c,\rho,w,d),
\end{equation}
where $s$ is its location in the atlas, $\tau$ its type, $c$ its scope and context, $\rho$ its evidential
reliability, $w$ its importance relative to $Q$, and $d$ the data, proof state, or argument supporting it.
\end{definition}

Important pressure types include:
\begin{description}[leftmargin=2.6cm,style=nextline]
  \item[Empirical pressure.] An anomaly, unexplained regularity, failed prediction, or measurement that
  separates otherwise plausible theories.
  \item[Inferential pressure.] A blocked proof, repeated lemma pattern, unresolved dependence, or
  computational bottleneck.
  \item[Structural pressure.] Noncommutation, failure of naturality, failed gluing, nonexactness, or an
  unrepresented functor.
  \item[Explanatory pressure.] Many independent laws or mechanisms account for related phenomena
  without explaining why they co-occur.
  \item[Unification pressure.] Parallel constructions, duplicated axioms, or recurring translations suggest
  a shared generative structure.
  \item[Representational pressure.] Results depend excessively on coordinates, presentation, gauge,
  encoding, or an unstable distinction between concepts.
  \item[Generative pressure.] A theory accommodates known cases but yields few new questions,
  predictions, or transferable methods.
  \item[Pragmatic pressure.] A decision, intervention, experiment, computation, or design task cannot be
  resolved efficiently by the present framework.
\end{description}

\subsection{Pressure as a context-indexed covector}

Let $\Def_D$ denote admissible local changes or completions of a diagram $D$. Conceptually, a pressure
acts like a covector: it evaluates how relevant a possible movement in theory space would be. In a
discrete implementation, this is simply a valuation
\begin{equation}
\pi_Q:\Def_D\longrightarrow V,
\end{equation}
where $V$ is a partially ordered space of relevance dimensions. In a geometric or derived setting,
$\Def_D$ may be replaced by a tangent complex or deformation object.

The pairing
\begin{equation}
\langle \pi_Q,h\rangle
\end{equation}
measures the expected relief or leverage of a candidate deformation $h$. This gives mathematical
content to the phrase ``follow the conceptual pressure'' without assuming that theory space is a
smooth manifold.

Alternatively, define a pressure functional
\begin{equation}
\Phi_Q(D)=\sum_{j}\lambda_j\phi_j(D),
\end{equation}
where the $\phi_j$ measure unresolved observations, blocked dependencies, duplicated assumptions,
representational sensitivity, and pragmatic cost. A completion $h:D\to D_h$ has pressure-relief vector
\begin{equation}
\Delta\Phi_Q(h)=\Phi_Q(D)-\Phi_Q(D_h).
\end{equation}
This quantity guides attention; it does not prove that $D_h$ is correct.

\subsection{From informal pressure to a formal defect}

The interpretation of pressure is itself generative. Scientific problems are usually stated in natural
language, equations, proof traces, simulation outputs, and expert disagreements rather than as labeled
categorical diagrams. Table~\ref{tab:pressure} gives typical translations.

\begin{table}[ht]
\centering
\small
\begin{tabularx}{\textwidth}{>{\raggedright\arraybackslash}p{0.29\textwidth}
>{\raggedright\arraybackslash}p{0.30\textwidth}X}
\toprule
Observed conceptual pressure & Candidate formal diagnosis & Generative questions \\
\midrule
The same ad hoc construction appears in unrelated proofs & Missing universal factor or representing object & What are all instances natural in? Through what minimal object do they factor? \\
Equivalent descriptions lead to different conclusions & Noncommuting square or missing higher coherence & Which path-dependent residue remains? Is it observable or representational? \\
Local models work but no global model does & Failed descent or nontrivial obstruction & What cocycle, topology, boundary datum, or anomaly measures the failure? \\
A parameter must be retuned independently at each scale & Missing naturality under scale change or incomplete effective variables & Is there an adjunction, flow, or latent invariant controlling the retuning? \\
Several laws explain correlated phenomena independently & Explanatory fragmentation & Can the laws be derived as specializations of a common mechanism with fewer independent assumptions? \\
A proof programme repeatedly stalls at the same interface & Dependency bottleneck or absent bridge theorem & What weakest lemma or interface semantics would close the largest downstream set? \\
A formalism is elegant but yields no new calculation & Missing concrete-return morphism & Which observable, theorem, algorithm, or experiment would distinguish it from a relabeling? \\
All completions inside the current vocabulary are fragile & Representational lock-in & Which object, morphism, identity criterion, or logic must be changed? \\
\bottomrule
\end{tabularx}
\caption{Interpreting informal conceptual pressures as formal search opportunities.}
\label{tab:pressure}
\end{table}

\section{A relevance calculus}

\subsection{Three gates before ranking}

Before assigning detailed scores, a candidate should pass three conjunctive gates:
\begin{equation}
\Rel_{\mathrm{gate}}(H\mid Q)
=
\Press_Q(H)\otimes
\mathsf{Leverage}_Q(H)\otimes
\mathsf{Return}_Q(H).
\end{equation}

Here $\otimes$ is a conjunction-like operation, such as a meet or geometric aggregation, chosen so
that a near-zero value in any factor cannot be concealed by a large value elsewhere.
\begin{itemize}
  \item $\Press_Q(H)$ asks whether $H$ addresses a material, well-supported pressure.
  \item $\mathsf{Leverage}_Q(H)$ asks whether resolving it affects important downstream structures.
  \item $\mathsf{Return}_Q(H)$ asks whether the abstraction can be pushed back to concrete
  consequences.
\end{itemize}

This prevents a highly connected but ungrounded abstraction from dominating the search merely because
it creates many formal relationships.

\subsection{A multi-objective relevance vector}

Candidates that pass the gates are compared using
\begin{equation}
\mathbf{R}_Q(H)=
\left(
R_{\mathrm{ub}},R_{\mathrm{exp}},R_{\mathrm{uni}},R_{\mathrm{pred}},
R_{\mathrm{trans}},R_{\mathrm{prag}},R_{\mathrm{nec}},R_{\mathrm{rob}}
\,;\,
-C_{\mathrm{comp}},-C_{\mathrm{frag}},-C_{\mathrm{aux}},-C_{\mathrm{risk}}
\right).
\label{eq:relevance-vector}
\end{equation}

The system should maintain a Pareto frontier rather than immediately collapsing
\eqref{eq:relevance-vector} to one scalar. Scalarization is useful only after the research context and
resource decision are explicit.

\subsection{Unblocking relevance}

Let $G_Q$ be a dependency hypergraph of open proof obligations, modeling tasks, experimental
decisions, and downstream goals. Let $B_Q$ denote blocked nodes and $D_H$ the obligations directly
resolved by $H$. Define
\begin{equation}
R_{\mathrm{ub}}(H)
=
\sum_{x\in \Cl_{G_Q}(D_H)\cap B_Q}w_Q(x)
-\lambda\Cost_Q(H).
\end{equation}

The closure $\Cl_{G_Q}(D_H)$ captures indirect effects: a modest bridge lemma may be more relevant
than a profound isolated theorem if it unlocks an entire proof chain. To avoid superficial gaming,
an obligation counts as resolved only after the relevant proof, computation, or empirical test has
actually passed.

\subsection{Explanatory power}

An explanation should do more than entail observations. It should reduce independent assumptions,
identify a mechanism or structural necessity, support counterfactual reasoning, and distinguish cases
that would otherwise be conflated. A schematic operator is
\begin{align}
R_{\mathrm{exp}}(H)
=\;&\sum_{e\in\mathcal{E}_Q}w_e\,
\mathsf{Derive}(H,e)
+\alpha\,\mathsf{Counterfactual}(H,e) \\
&-\beta\,\mathsf{AuxiliaryAssumptions}(H,e)
-\gamma\,\mathsf{Overfit}(H,e).
\end{align}

Explanatory compression resembles accounts of unification in which many phenomena are derived from a
smaller set of argument patterns \cite{kitcher1981}, but compression alone is insufficient. A physical
explanation must retain an operational connection to observables, and a mathematical explanation
should illuminate why the result follows rather than merely shorten a proof term.

\subsection{Unification power}

Unification power is high when independent theories or derivations become specializations of a common
structure and when nontrivial results can be transported through it. Define
\begin{align}
R_{\mathrm{uni}}(H)
=\;&a\,\Delta\mathsf{DescriptionLength}(H)
+b\,\mathsf{SharedFactorization}(H) \\
&+c\,\mathsf{CrossDomainTransport}(H)
-d\,\mathsf{ForcedIdentifications}(H).
\end{align}

The last term penalizes unification obtained by erasing distinctions that matter empirically or
mathematically. High unification value requires recoverability: the original theories should reappear as
controlled specializations, images, limits, or sectors of the proposed framework.

\subsection{Predictive and generative power}

Let $\mathsf{Cons}(H)$ be the consequences generated from $H$ and let $\Cl(E)$ be the deductive or
statistical closure of evidence used to construct it. Then
\begin{equation}
R_{\mathrm{pred}}(H)
=\sum_{q\in\mathsf{Cons}(H)\setminus\Cl(E)}
w_q\,\mathsf{Independence}(q,E)\,
\mathsf{Discrimination}(q)\,
\mathsf{Testability}(q).
\end{equation}

For pure mathematics, a prediction can be a new conjecture, a reusable lemma, a classification, a
counterexample, or an algorithm. For physics, it should ideally identify a quantitative regime in which
the proposal differs from live alternatives.

\subsection{Theorem-transport relevance}

A concept can be valuable because it turns a proof in one domain into a family of results elsewhere.
If $F_i$ are validated theory morphisms, define
\begin{equation}
R_{\mathrm{trans}}(H)
=\sum_i\sum_{\varphi\in\mathcal{T}_H}
w_{i,\varphi}\,
\mathbf{1}\!\left[
F_i(\varphi)\text{ is well-typed, nontrivial, and newly derivable}
\right].
\end{equation}

Transport must preserve meaning and scope, not merely syntax. The satisfaction condition from
institution theory is a natural formal requirement for this layer \cite{goguenburstall1992}.

\subsection{Pragmatic relevance}

Pragmatic value includes new experiments, improved decisions, reduced computation, better model
selection, or feasible interventions. A decision-theoretic form is
\begin{equation}
R_{\mathrm{prag}}(H)
=
\max_a\mathbb{E}[U(a)\mid H,E]
-\max_a\mathbb{E}[U(a)\mid E]
-\Cost_{\mathrm{use}}(H).
\end{equation}

This operator is intentionally context-specific. It must not be confused with epistemic truth: an
instrumentally useful effective model can have high pragmatic relevance without being a fundamental
ontology.

\subsection{Necessity, robustness, and anti-sterility}

Conceptual necessity is high when a candidate is selected by a universal property, a minimal
obstruction, a uniqueness theorem, or a conservative minimality result. Robustness asks whether the
candidate survives equivalent encodings, reasonable changes of granularity, and perturbations of the
evidential set.

An anti-sterility condition requires at least one concrete return path:
\begin{enumerate}
  \item a blocked proof obligation is resolved;
  \item a new discriminating prediction or conjecture is produced;
  \item an existing result is derived from strictly weaker or fewer assumptions;
  \item a simulation, experiment, or computation becomes possible;
  \item a useful intervention or decision changes;
  \item a counterexample or no-go result eliminates a substantial search region.
\end{enumerate}

The marginal contribution of $H$ should also be tested by ablation:
\begin{equation}
R_{\mathrm{abl}}(H)
=\Val_Q(D+H)-\Val_Q(D),
\end{equation}
with the rest of the theory and evidence held fixed as far as possible.

\section{The relevance-guided generative element}

\subsection{A constrained proposal kernel}

The generative element should not be modeled as a theorem prover or a deterministic categorical
operator. It is a context-sensitive proposal mechanism
\begin{equation}
\Gen_\theta
\bigl(H\mid D,\mathcal{P}_Q,Q,\mathcal{K},\mathcal{H}\bigr),
\end{equation}
where $D$ is the current atlas region, $\mathcal{P}_Q$ the active pressure field, $\mathcal{K}$ the
allowed completion grammar, and $\mathcal{H}$ the history of attempts, failures, and validations.
Its output is a distribution or ranked portfolio over typed hypothesis extensions.

The generator has four modes:
\begin{enumerate}
  \item \textbf{Local completion:} fill a missing object, morphism, law, or proof interface inside the
  present representation.
  \item \textbf{Bridge induction:} propose a functor, adjunction, duality, shared interface, or common
  universal property between previously separate theories.
  \item \textbf{Mechanism induction:} introduce a latent state, interaction, invariant, or dynamical law
  that explains several pressures jointly.
  \item \textbf{Representation mutation:} alter the object language, morphism types, identity criterion,
  scope modalities, or foundational logic when the existing atlas systematically produces fragile
  completions.
\end{enumerate}

Only the first mode resembles ordinary completion. The fourth is necessary for nonincremental
conceptual change.

\subsection{The hypothesis package}

Every generated candidate should be compiled into a typed package
\begin{equation}
H=(s,p,m,\widehat D,\eta,\mathcal{C},\mathcal{F},\kappa,u),
\end{equation}
containing:
\begin{itemize}
  \item the pressure site $s$ and pressure items $p$ it claims to address;
  \item the formal move $m$ and extended diagram $\widehat D$;
  \item an informal explanatory rationale $\eta$;
  \item expected consequences $\mathcal{C}$;
  \item explicit falsifiers, counterexamples, or failure conditions $\mathcal{F}$;
  \item implementation and validation cost $\kappa$;
  \item uncertainty $u$, including uncertainty about the formalization itself.
\end{itemize}

This prevents the generator from producing an attractive concept with no account of what it changes,
why it is relevant, or how it could fail.

\subsection{Selection without a single reward}

A single learned reward for ``interesting theory'' would be vulnerable to Goodhart effects. It would
favor fluent, highly connected, fashionable, or excessively general proposals. The controller should
instead:
\begin{enumerate}
  \item apply the pressure--leverage--return gates;
  \item enforce formal and empirical admissibility thresholds;
  \item retain a Pareto frontier over the relevance vector;
  \item allocate validation effort by expected value of information;
  \item reserve an exploration budget for candidates outside the current dominant representation;
  \item preserve negative results and failed candidates so that the same sterile abstractions are not
  repeatedly regenerated.
\end{enumerate}

A stochastic portfolio is desirable. Scientific relevance is uncertain, and premature convergence on
one elegant unification can suppress alternative representations that later prove essential.

\section{Philosophical reflection at essential fragmentation points}

\subsection{Meaning before mechanics}

The preceding machinery still risks beginning too late. A noncommuting diagram is visible only after
the relevant objects and morphisms have been chosen. Yet the deepest conceptual pressure often concerns
the prior choice of what counts as an object, an identity, a state, an observable, an explanation, or a
legitimate transformation. At such points the reasoner must temporarily suspend the completion grammar
and ask what the competing concepts mean in their native theoretical practices.

\begin{principle}[Semantic priority]
When a pressure concerns the identity, explanatory role, or operational interpretation of a concept,
the system must reconstruct that concept's meaning profile before selecting a mathematical completion
operator.
\end{principle}

This does not place informal reflection outside the architecture. It gives reflection a typed output
which can subsequently constrain formalization. The philosophical layer is generative because it can
redescribe the problem, alter the primitives, and expose a pressure that was invisible in the original
atlas.

\subsection{Meaning profiles}

The meaning of a scientific concept cannot be represented adequately by a textual definition or a
vector embedding. It is distributed across its inferential, referential, operational, modal, and
pragmatic roles. For a concept $c$ in theory fragment $T$ and research context $Q$, define
\begin{equation}
\mu_Q(c,T)=
\bigl(
\mathsf{Ref},\mathsf{Infer},\mathsf{Id},\mathsf{Modal},
\mathsf{Op},\mathsf{Scale},\mathsf{Explain},
\mathsf{Use},\mathsf{Exclude}
\bigr).
\label{eq:meaning-profile}
\end{equation}
The components record:
\begin{itemize}
  \item what the concept purports to refer to, if anything;
  \item which inferences license its introduction and which follow from it;
  \item the criterion by which two instances count as the same;
  \item its counterfactual and modal commitments;
  \item its measurement, construction, or computational anchors;
  \item the scales and regimes in which it is meaningful;
  \item the explanatory role it performs;
  \item the interventions, predictions, or mathematical operations it enables;
  \item the distinctions it is intended to preserve or rule out.
\end{itemize}

This plural profile avoids committing the system to one philosophical theory of meaning. Referential,
inferentialist, operational, structural, and pragmatic aspects are retained as distinguishable
constraints. The inclusion of inferential commitment and relational function draws on traditions
exemplified by Brandom and Cassirer without reducing scientific meaning to either account
\cite{brandom1994,cassirer1910}. Their disagreement is often itself the source of discovery. For
example, two formulations
may agree extensionally on all current observations while assigning different identity criteria or
counterfactual roles to the same term. The disagreement is scientifically dormant until a new regime
or intervention makes it visible.

\subsection{Fragmentation and essential fragmentation}

Let $c_i$ be local realizations of an apparent common concept across contexts $T_i$, with partial
transport maps
\begin{equation}
\tau_{ij}:\mu_Q(c_i,T_i)\dashrightarrow\mu_Q(c_j,T_j).
\end{equation}
A fragmentation site occurs when these transports are missing, path dependent, scope violating, or
unable to preserve a protected component of the meaning profile.

Fragmentation has several forms:
\begin{description}[leftmargin=3.0cm,style=nextline]
  \item[Lexical fragmentation.] Different vocabularies express compatible structures but no reliable
  translation has been articulated.
  \item[Representational fragmentation.] The same content is encoded through nonaligned state spaces,
  variables, gauges, or mathematical foundations.
  \item[Inferential fragmentation.] Results cannot be transported even though the concepts appear to
  refer to the same objects.
  \item[Explanatory fragmentation.] Theories agree on predictions but disagree about why they hold or
  which counterfactuals sustain them.
  \item[Ontological fragmentation.] The theories disagree about what is fundamental, derivative, or
  physically real.
  \item[Operational fragmentation.] A concept has incompatible measurement, construction, or
  intervention procedures.
  \item[Scale fragmentation.] The identity and role of a concept change across regimes without an
  explicit transition law.
  \item[Global fragmentation.] Locally coherent meanings fail to glue into a global interpretation.
\end{description}

Not every fragmentation is conceptually important. Some disappear under renaming, unit conversion,
definitional expansion, gauge transformation, conservative extension, or a known equivalence.
Let $\mathcal{R}_{\mathrm{cheap}}$ be a declared class of such low-cost repairs and
$\delta_Q(D,s)$ a vector of semantic and inferential mismatch at site $s$. Define the residual
fragmentation
\begin{equation}
\mathsf{EssFrag}_Q(s)
=
\inf_{r\in\mathcal{R}_{\mathrm{cheap}}}
\left\lVert \delta_Q(rD,s)\right\rVert.
\label{eq:essential-fragmentation}
\end{equation}
A fragmentation is \emph{essential relative to $Q$} when this residual remains high while the site
also has substantial downstream consequence:
\begin{align}
\mathsf{FragValue}_Q(s)
=\;&\mathsf{EssFrag}_Q(s)
\otimes\mathsf{Centrality}_Q(s) \\
&\otimes\mathsf{ConsequenceDivergence}_Q(s)
\otimes\mathsf{ConcreteReach}_Q(s).
\end{align}

The essentiality judgment must identify what would be lost by a proposed repair. If two ontologies can
be made formally equivalent only by discarding their intervention semantics, the fragmentation has not
been dissolved; it has been hidden.

\subsection{Philosophical reflection lenses}

Deep reflection should not be one unrestricted prompt asking a model to think philosophically.
The system applies several lenses independently and then compares their outputs.

\begin{table}[ht]
\centering
\small
\begin{tabularx}{\textwidth}{
>{\raggedright\arraybackslash}p{0.19\textwidth}
>{\raggedright\arraybackslash}p{0.33\textwidth}X}
\toprule
Lens & Central question & Required output \\
\midrule
Hermeneutic & How does the part acquire meaning from the whole theory, and how does it revise that whole? & Context-dependent meanings, translation losses, and alternative readings \\
Phenomenological--operational & What is actually given, measured, constructed, or experienced before theoretical interpretation? & Operational anchors and a separation of datum from interpretation \\
Constitutive & What must already be presupposed for the theory's objects, measurements, or explanations to be possible? & Hidden background structures and conditions of intelligibility \\
Dialectical--immanent & Which commitments of the theory generate internal tension when jointly developed? & Minimal contradiction or tension and possible transformations that preserve both insights \\
Genealogical & Which historical problems and inherited distinctions produced the present conceptual division? & Contingent versus necessary distinctions and plausible alternative developments \\
Pragmatic & What difference would adopting the concept make to inquiry, intervention, calculation, or decision? & Actionable consequences and cases in which the distinction is idle \\
Structural realist & Which relations or invariants survive changes of representation and ontology? & Candidate preserved structure and the maps under which it is invariant \\
Modal--counterfactual & In which possible or counterfactual cases do the competing concepts separate? & Discriminating counterworlds, boundary cases, and modal commitments \\
\bottomrule
\end{tabularx}
\caption{Reflection lenses used to interpret essential conceptual fragmentation.}
\label{tab:reflection-lenses}
\end{table}

These lenses need not converge. Their disagreement becomes a structured pressure. A pragmatic
equivalence combined with an ontological divergence, for example, calls for a different research move
than an inferential equivalence combined with an operational divergence. Philosophical pluralism is
therefore represented as a portfolio of constrained views rather than averaged into a bland synthesis.
The use of historical and immanent reconstruction also guards against treating the current formal
vocabulary as inevitable, a concern emphasized in different ways by Kuhn and Lakatos
\cite{kuhn1962,lakatos1976}.

\subsection{Generative reflection operators}

The reflection layer acts through operators that can change the conceptual organization before a
categorical construction is chosen:
\begin{description}[leftmargin=3.2cm,style=nextline]
  \item[Ontological suspension.] \(\mathsf{Suspend}\) temporarily withholds the dominant ontology and restates the pressure in
  terms of observations, inferential roles, and practices.
  \item[Primitive shift.] \(\mathsf{PrimitiveShift}\) promotes a relation, process, transformation, or constraint to a
  primitive, or reconstruct an apparent object as relationally generated.
  \item[Dependency inversion.] \(\mathsf{DependencyInvert}\) exchanges explanans and explanandum: ask whether what was treated
  as a background law is generated by the states it was supposed to govern.
  \item[Context internalization.] \(\mathsf{ContextInternalize}\) turns an implicit background, scale, observer, coordinate
  system, or experimental arrangement into an explicit variable or indexed context.
  \item[Level shift.] \(\mathsf{LevelShift}\) moves between object, process, law, constraint on laws, and
  transformation between theories.
  \item[Identity revision.] \(\mathsf{IdentityRevise}\) alters the criterion of sameness: strict equality may become
  equivalence, path-connectedness, observational indistinguishability, or coherent transport.
  \item[Invariant extraction.] \(\mathsf{InvariantExtract}\) searches for what survives the competing representations rather
  than choosing one representation as fundamental.
  \item[Obstruction reification.] \(\mathsf{ObstructionReify}\) treats a failure of translation or gluing as a positive object
  carrying information rather than as noise to eliminate.
  \item[Dualization.] \(\mathsf{Dualize}\) exchanges state and observable, syntax and semantics, local and global,
  kinematics and dynamics, or construction and constraint.
  \item[Boundary probing.] \(\mathsf{BoundaryProbe}\) moves to singular limits, extreme scales, degenerate examples, and
  interface cases in which implicit assumptions become visible.
  \item[Counterworld construction.] \(\mathsf{Counterworld}\) constructs the smallest coherent model in which two apparently
  equivalent meanings give different consequences.
  \item[Operational reanchoring.] \(\mathsf{OperationalReanchor}\) requires a proposed redescription to change or simplify a
  measurement, proof, calculation, intervention, or experimental discrimination.
\end{description}

The operators can be composed, but the composition must remain traceable. A proposed theory may arise,
for example, by internalizing the descriptive context, revising identity from equality to coherent
transport, and then extracting the invariant of that transport. This is a materially different
discovery path from directly taking a pushout of the original theories.

\subsection{Compiling reflection into formal transformations}

Reflection becomes part of a rigorous reasoner only when its outputs generate formal obligations.
Table~\ref{tab:reflection-compilation} gives a nonexclusive compilation map.

\begin{table}[ht]
\centering
\small
\begin{tabularx}{\textwidth}{
>{\raggedright\arraybackslash}p{0.25\textwidth}
>{\raggedright\arraybackslash}p{0.32\textwidth}X}
\toprule
Reflective move & Candidate formal expression & Immediate obligation \\
\midrule
Internalize context & Indexed category, fibration, modality, dependent type & Show how old context-specific claims arise by reindexing \\
Revise identity & Localization, quotient, groupoid completion, univalence & Prove which distinctions are preserved and which are intentionally identified \\
Promote relations to primitives & Enriched category, profunctor, process theory, operad & Reconstruct former objects and their observables \\
Extract an invariant & Limit, end, natural transformation, representable functor & Establish invariance and nontriviality under validated transformations \\
Reify failed gluing & Homotopy fiber, obstruction class, cohomology, descent datum & Compute the obstruction in at least one concrete model \\
Invert dependence & Adjunction, duality, fixed-point problem, coupled dynamics & Derive both directions and prevent circular explanation \\
Shift theoretical level & Internal language, meta-theory morphism, reflective tower & State type levels and block self-referential collapse \\
Construct a counterworld & Model finder, free model, forcing or perturbation family & Verify consistency and isolate the discriminating commitment \\
\bottomrule
\end{tabularx}
\caption{From philosophical redescription to formal operator and proof obligation.}
\label{tab:reflection-compilation}
\end{table}

The compilation relation is many-to-many. No philosophical lens mechanically determines a category
theoretic construction. Instead, it restricts which constructions preserve the meaning at stake and
which proof obligations must be discharged.

\subsection{The deep-reflection cycle}

At an essential fragmentation point, the system executes the following cycle:
\begin{enumerate}
  \item \textbf{Native reconstruction.} Restate each theory in its own terms, with sources and scope.
  \item \textbf{Semantic profiling.} Construct the meaning profiles in
  \eqref{eq:meaning-profile}.
  \item \textbf{Cheap-repair test.} Attempt definitional, representational, and conservative repairs.
  \item \textbf{Essentiality localization.} Record the residual mismatch and its downstream importance.
  \item \textbf{Plural reflection.} Apply the lenses in Table~\ref{tab:reflection-lenses} independently.
  \item \textbf{Immanent challenge.} Let each lens challenge the assumptions and losses of the others.
  \item \textbf{Redescription portfolio.} Generate several conceptual reorganizations using the
  reflection operators.
  \item \textbf{Counterworld stress.} Construct cases in which the redescriptions make divergent
  commitments.
  \item \textbf{Formal compilation.} Translate each surviving redescription into formal objects,
  morphisms, and obligations.
  \item \textbf{Concrete return.} Demand a proof, model, computation, prediction, or intervention that
  would change if the redescription were adopted.
\end{enumerate}

The cycle is hermeneutic in a precise operational sense: local concepts are interpreted through the
whole theory, candidate revisions alter the whole, and the altered whole changes the local meaning.
Iteration stops when further passes do not materially change the protected meaning profiles, active
pressure set, candidate frontier, or concrete discriminations.

\subsection{The reflection record}

Each reflection pass produces
\begin{equation}
\mathcal{J}=
(s,\mu,\delta,L,\rho,\Pi,\Omega,\mathcal{O},u),
\end{equation}
where $s$ is the fragmentation site, $\mu$ the meaning profiles, $\delta$ the residual mismatch, $L$
the lens outputs, $\rho$ the proposed redescriptions, $\Pi$ the protected commitments, $\Omega$ the
anticipated losses, $\mathcal{O}$ the formalization and concrete-return obligations, and $u$ the
uncertainty.

A reflection record is rejected as shallow when it:
\begin{itemize}
  \item merely restates the source terminology;
  \item names a familiar philosophical position without deriving consequences;
  \item proposes a synthesis by deleting the disputed distinction;
  \item lacks a counterworld in which alternatives separate;
  \item cannot identify what the new representation preserves and sacrifices;
  \item has no route to a formal or concrete obligation.
\end{itemize}

This is the main defense against decorative philosophy. Reflection earns a role in discovery by
changing the admissible transformations, the ranking of candidates, or the concrete tests.

\section{Agentic implementation}

\subsection{Implementation objective and design constraints}

The implementation target is a local-first, provider-neutral research system that turns a scoped body
of mathematical or physical material into a validated portfolio of conceptual hypotheses. It must
support open-ended generation without allowing generated prose to alter the canonical theory atlas
directly.

The following constraints are architectural:
\begin{enumerate}
  \item \textbf{Typed boundaries.} Every skill has versioned input and output schemas.
  \item \textbf{Separation of epistemic roles.} Generators propose; formal, computational, and empirical
  components validate.
  \item \textbf{Provenance.} Every extracted claim, interpretation, and formal object points to its
  source or is explicitly labeled as conjectural.
  \item \textbf{Scope preservation.} Regime and approximation metadata travel with every claim.
  \item \textbf{Plurality.} Competing interpretations remain separate until an accepted equivalence or
  synthesis is established.
  \item \textbf{Replayability.} Model, prompt, skill version, parameters, inputs, and random seed are
  recorded for every generative action.
  \item \textbf{Budgeted persistence.} Every loop has termination conditions, resource limits, and an
  explicit unresolved outcome.
  \item \textbf{No graph-only semantics.} The graph indexes relationships, while typed declarations,
  proofs, models, evidence, and reflection records carry the meaning.
\end{enumerate}

\subsection{Logical architecture}

The system is organized into six planes:
\begin{description}[leftmargin=3.0cm,style=nextline]
  \item[Knowledge plane.] Versioned theory fragments, morphisms, higher cells, meaning profiles,
  evidence, examples, proof obligations, and source documents.
  \item[Interpretive plane.] Pressure items, fragmentation sites, philosophical lens outputs, protected
  commitments, and conceptual redescriptions.
  \item[Generative plane.] Candidate completions, bridges, mechanisms, counterworlds, and representation
  mutations.
  \item[Formal plane.] Type checking, theorem proving, model finding, categorical computations, and
  transport verification.
  \item[Concrete plane.] Symbolic calculation, numerical simulation, benchmark instances, experimental
  design, and pragmatic evaluation.
  \item[Control plane.] Workflow state, budgets, Pareto selection, uncertainty, drift monitoring,
  human review, and final synthesis.
\end{description}

The planes interact through immutable artifacts and typed references. A free-text reflection can be
stored as part of a reflection record, but only its explicitly extracted commitments, losses,
obligations, and candidate transformations may enter downstream formal workflows.

\subsection{Agentic functions and authority boundaries}

The system uses the following agentic functions:
\begin{description}[leftmargin=3.3cm,style=nextline]
  \item[Atlas cartographer.] Extracts definitions, claims, models, observables, scopes, examples, and
  proposed translations. It may stage additions but cannot approve merges.
  \item[Pressure interpreter.] Converts anomalies, proof traces, expert questions, and pragmatic needs
  into weighted pressure items.
  \item[Meaning profiler.] Reconstructs the meaning profile of contested concepts in each native
  theoretical context.
  \item[Fragmentation analyst.] Performs cheap-repair tests and identifies residual, high-consequence
  fragmentation.
  \item[Philosophical reflector.] Runs the independent lenses, generates redescriptions, states losses,
  and constructs discriminating counterworlds.
  \item[Structural diagnostician.] Compiles reflections into representability, coherence, descent,
  factorization, adjunction, obstruction, or meta-theory problems.
  \item[Generative theorist.] Produces a diverse portfolio of typed hypothesis packages across local,
  bridge, mechanism, and representation-mutation modes.
  \item[Relevance analyst.] Applies the gates and computes the multi-objective relevance profile without
  making truth judgments.
  \item[Formalizer and prover.] Checks typing, universal properties, conservative extension, consistency
  relative to the selected foundation, and theorem transport.
  \item[Concrete instantiator.] Derives equations, examples, simulations, observables, algorithms, or
  experiments from the abstract proposal.
  \item[Adversarial critic.] Searches for countermodels, vacuity, scope violations, hidden auxiliary
  assumptions, equivalent prior art, and rival explanations.
  \item[Meta-representation critic.] Determines whether repeated failure is caused by the current
  ontology or morphism grammar and proposes a new atlas schema.
  \item[Portfolio controller.] Allocates effort, maintains diversity, requests human judgment, and
  promotes only candidates that satisfy declared thresholds.
\end{description}

No function validates its own generative output. The portfolio controller may schedule tasks and apply
declared policy, but it may not override a failed proof, falsifying model, or absent concrete return
path.

\subsection{Canonical data contracts}

The canonical objects should be represented as versioned, schema-validated records. The following
Python-like definitions are indicative rather than tied to one library:

\begin{lstlisting}[style=spec,language=Python,caption={Core typed records.}]
class EvidenceRef:
    artifact_id: str
    locator: str
    evidence_type: Literal[
        "proof", "paper", "data", "simulation", "expert", "conjecture"
    ]
    reliability: float

class MeaningProfile:
    concept_id: str
    theory_id: str
    referential_role: list[str]
    inferential_commitments: list[str]
    identity_criteria: list[str]
    modal_commitments: list[str]
    operational_anchors: list[str]
    scale_and_scope: list[str]
    explanatory_role: list[str]
    pragmatic_uses: list[str]
    protected_distinctions: list[str]
    evidence: list[EvidenceRef]
    uncertainty: dict[str, float]

class FragmentationSite:
    site_id: str
    concept_ids: list[str]
    theory_ids: list[str]
    mismatch_vector: dict[str, float]
    cheap_repairs_attempted: list[str]
    residual_fragmentation: float
    protected_losses: list[str]
    downstream_targets: list[str]

class ReflectionRecord:
    site_id: str
    lens_outputs: dict[str, str]
    redescriptions: list[str]
    preserved_commitments: list[str]
    anticipated_losses: list[str]
    counterworlds: list[str]
    formal_obligations: list[str]
    concrete_obligations: list[str]
    source_refs: list[EvidenceRef]

class HypothesisPackage:
    candidate_id: str
    pressure_ids: list[str]
    reflection_ids: list[str]
    formal_move: str
    atlas_patch: dict
    consequences: list[str]
    falsifiers: list[str]
    relevance_vector: dict[str, float]
    validation_tasks: list[str]
    uncertainty: dict[str, float]
\end{lstlisting}

Additional first-class records represent theory fragments, theory morphisms, higher cells, research
contexts, pressure items, validation outcomes, task dependencies, and run states. Every record has a
stable identifier, schema version, creation provenance, and parent versions. Its lifecycle status is
one of
\begin{equation}
\mathsf{staged},\qquad
\mathsf{accepted},\qquad
\mathsf{rejected},\qquad
\mathsf{superseded}.
\end{equation}

\subsection{Skill contract}

A skill is a bounded epistemic operation, not merely a prompt. Its specification is
\begin{equation}
\mathcal{S}=
(\mathsf{name},\mathsf{version},I,O,P,Q,T,B,E),
\end{equation}
where $I$ and $O$ are input and output schemas, $P$ the preconditions, $Q$ the postconditions, $T$
the authorized tools, $B$ the resource budget, and $E$ the termination and escalation policy.

Every skill directory should contain:
\begin{lstlisting}[style=spec,caption={Required skill directory contract.}]
skills/<skill-name>/
  SKILL.md                 # purpose, method, boundaries
  input.schema.json
  output.schema.json
  policy.yaml              # tools, budget, stop and escalation rules
  prompts/
    system.md
    task.md
    critique.md
  tests/
    fixtures/
    test_contract.py
    test_invariance.py
    test_adversarial.py
  examples/
    minimal_input.json
    expected_output.json
\end{lstlisting}

All skill results return a common envelope:
\begin{equation}
\mathsf{SkillResult}=
(\mathsf{status},\mathsf{artifacts},\mathsf{claims},
\mathsf{uncertainties},\mathsf{next},\mathsf{provenance}).
\end{equation}
The status vocabulary includes \(\mathsf{completed}\), \(\mathsf{partial}\),
\(\mathsf{blocked}\), and \(\mathsf{invalid}\). A blocked result is a valid terminal outcome and must
state the missing information or authority.

\subsection{Initial skill catalogue}

\begin{longtable}{@{}
>{\raggedright\arraybackslash}p{0.23\textwidth}
>{\raggedright\arraybackslash}p{0.31\textwidth}
>{\raggedright\arraybackslash}p{0.38\textwidth}@{}}
\toprule
Skill & Primary input and output & Core validation \\
\midrule
\endfirsthead
\toprule
Skill & Primary input and output & Core validation \\
\midrule
\endhead
Scope research context & Question, corpus, goals, budgets \(\to Q\) & Scope is explicit; success and stop criteria are measurable \\
Map theory fragment & Source artifacts \(\to\) typed theory fragments & Every nonconjectural claim has a source; types and regimes are retained \\
Propose theory morphism & Two fragments \(\to\) staged interpretation map & Symbol mapping is total on declared domain; satisfaction obligations emitted \\
Interpret conceptual pressure & Goals, anomalies, proof states \(\to\) pressure items & Pressure is localized, evidenced, and distinguished from preferred solution \\
Profile concept meaning & Concept plus native sources \(\to\) meaning profile & All nine components of \eqref{eq:meaning-profile} are filled or marked unknown \\
Detect essential fragmentation & Meaning profiles and translations \(\to\) fragmentation sites & Cheap repairs attempted; protected losses and residual mismatch reported \\
Run philosophical reflection & Fragmentation site \(\to\) reflection records & Independent lenses, counterworld, losses, and obligations are present \\
Diagnose categorical defect & Atlas diagram and reflection \(\to\) formal defect candidates & Required objects and morphisms are typed; metaphorical mappings are flagged \\
Generate theory candidates & Pressures, reflections, defects \(\to\) hypothesis portfolio & Minimum diversity across generation modes; duplicates quotient by equivalence \\
Evaluate relevance & Candidate, context, dependency graph \(\to\) relevance vector & Gates reported separately; no truth claim inferred from score \\
Formalize candidate & Hypothesis package \(\to\) formal declarations and obligations & Type checking passes or exact blockers are returned \\
Prove transport & Morphism and theorem set \(\to\) transport results & Scope and satisfaction preserved; trivial restatements excluded \\
Instantiate concretely & Formal candidate and target domains \(\to\) models or experiments & Return diagrams and error bounds supplied \\
Falsify candidate & Candidate and adversarial budget \(\to\) countermodels and stress report & Searches scope boundaries, degenerate cases, and rival explanations \\
Mutate representation & Repeated structured failures \(\to\) atlas-schema proposal & Existing content migration and information-loss analysis included \\
Synthesize portfolio & Validations and relevance profiles \(\to\) research recommendation & Pareto frontier, disagreement, uncertainty, and next discriminating actions retained \\
Audit reasoning run & Complete run trace \(\to\) integrity report & Provenance, role separation, budget, drift, and replay checks pass \\
\bottomrule
\caption{Initial skill catalogue for the theory-discovery system.}
\label{tab:skills}
\end{longtable}

\subsection{Critical skill procedures}

\paragraph{The philosophical-reflection skill.}
This skill runs one branch per lens in Table~\ref{tab:reflection-lenses}. Each branch receives the same
source-grounded meaning profiles but does not see the other branches' synthesis. A comparison pass then
identifies convergence, contradiction, and complementary insight. A final generative pass proposes at
least three redescriptions, including one that changes the primitive ontology and one conservative
alternative. The output fails validation if it contains no protected loss, counterworld, or downstream
obligation.

\paragraph{The essential-fragmentation skill.}
This skill runs repair attempts in increasing order of cost: lexical normalization, unit and notation
alignment, definitional extension, known equivalence, conservative interpretation, and scope
restriction. It estimates the residual only after these attempts. Human review is required when a
repair would change operational or ontological commitments.

\paragraph{The candidate-generation skill.}
Generation occurs in two passes. An exploratory pass samples diverse conceptual moves with higher
temperature or equivalent diversity controls. A normalization pass converts them into hypothesis
packages, merges equivalent candidates, and rejects packages without formalization hooks. The
generator may cite source claims or label novel conjectures; it may not blur the two.

\paragraph{The representation-mutation skill.}
Mutation is triggered by a pattern of failures rather than one unsuccessful proof. Its proposal must
specify the new object and morphism grammar, a migration functor from the current atlas, information
lost or newly expressible, and at least one previously sterile pressure that becomes tractable.

\subsection{Workflow state machine}

A research run advances through the following states:
\begin{equation}
\begin{split}
\mathsf{SCOPED}&\to\mathsf{MAPPED}\to\mathsf{PRESSURED}\to
\mathsf{FRAGMENTED}\to\mathsf{REFLECTED}\\
&\to\mathsf{GENERATED}\to\mathsf{GATED}\to
\mathsf{FORMALIZED}\to\mathsf{INSTANTIATED}\\
&\to\mathsf{TESTED}\to\mathsf{SYNTHESIZED}.
\end{split}
\label{eq:state-machine}
\end{equation}

Transitions are guarded:
\begin{itemize}
  \item \(\mathsf{MAPPED}\to\mathsf{PRESSURED}\) requires a versioned atlas region and research context.
  \item \(\mathsf{FRAGMENTED}\to\mathsf{REFLECTED}\) requires a passed essentiality or high-impact
  uncertainty threshold.
  \item \(\mathsf{GATED}\to\mathsf{FORMALIZED}\) requires pressure, leverage, and return gates.
  \item \(\mathsf{FORMALIZED}\to\mathsf{INSTANTIATED}\) permits partial formalization only when exact
  residual obligations are carried forward.
  \item \(\mathsf{TESTED}\to\mathsf{SYNTHESIZED}\) requires both supportive and adversarial results.
\end{itemize}

Failure transitions are productive. A failed formalization returns to \(\mathsf{REFLECTED}\) when the
concept is ill typed, to \(\mathsf{GENERATED}\) when only one candidate is defective, and to a
\(\mathsf{REPRESENTATION\_REVIEW}\) branch when several candidates fail for the same structural
reason. A falsified candidate is retained with its failure conditions and may delimit a no-go region.

\subsection{End-to-end workflows}

\paragraph{Workflow A: atlas construction.}
\begin{enumerate}
  \item Ingest papers, formal libraries, code, data descriptions, and expert notes.
  \item Segment them into source-addressable claims, definitions, equations, examples, and scopes.
  \item Build staged theory fragments and identify unresolved term alignments.
  \item Propose morphisms and generate satisfaction, typing, and scope obligations.
  \item Approve only validated fragments and maps into the canonical atlas revision.
\end{enumerate}

\paragraph{Workflow B: pressure and reflection.}
\begin{enumerate}
  \item Construct the task dependency graph and collect anomalies, disagreements, and pragmatic goals.
  \item Localize pressure items and estimate evidential reliability.
  \item Profile the meaning of concepts touching high-impact sites.
  \item Run cheap-repair and essential-fragmentation analysis.
  \item Apply plural philosophical lenses and generate reflection records.
  \item Compile the records into formal defects and redescription obligations.
\end{enumerate}

\paragraph{Workflow C: theory generation.}
\begin{enumerate}
  \item Generate candidates independently in local, bridge, mechanism, and mutation modes.
  \item Normalize candidates into typed packages and quotient obvious equivalences.
  \item Apply pressure--leverage--return gates.
  \item Compute the relevance vector and retain a diverse Pareto frontier.
  \item Schedule the cheapest discriminating validation for each frontier candidate.
\end{enumerate}

\paragraph{Workflow D: formal and concrete validation.}
\begin{enumerate}
  \item Compile the candidate into MMT, Lean, a generalized algebraic theory, or another declared
  formal target.
  \item Check type correctness, model existence, conservativity, universal properties, and theorem
  transport.
  \item Instantiate it in at least one mathematical and, where applicable, physical model.
  \item Derive held-out consequences and construct countermodels or boundary cases.
  \item Update the relevance profile using realized rather than anticipated value.
\end{enumerate}

\paragraph{Workflow E: synthesis or representation revision.}
The controller compares surviving candidates, unresolved disagreements, and validation cost. It emits
one of four outcomes: promote a candidate; continue a discriminating experiment; suspend with explicit
unknowns; or revise the atlas representation. A synthesis never hides a failed validation or collapses
non-equivalent interpretations.

\subsection{Concrete software architecture}

The initial implementation can be a modular monorepository:
\begin{lstlisting}[style=spec,caption={Proposed repository layout.}]
theory-discovery/
  pyproject.toml
  src/tdr/
    domain/              # canonical schemas and invariants
    atlas/               # fragments, morphisms, higher cells
    pressures/           # pressure and dependency analysis
    reflection/          # lenses and redescription operators
    generation/          # provider-neutral proposal kernel
    relevance/           # gates, vectors, Pareto selection
    orchestration/       # state machine, budgets, policies
    validation/          # formal, numerical, empirical adapters
    provenance/          # artifacts, source anchors, replay
    api/                 # HTTP and CLI interfaces
  skills/                # versioned skill directories
  adapters/
    mmt/
    lean/
    gatlab/
    cubical_agda/
    smt/
    simulation/
    language_models/
  schemas/
  benchmarks/
  tests/
  artifacts/             # local content-addressed development store
\end{lstlisting}

The recommended runtime components are:
\begin{itemize}
  \item a Python orchestration and API layer with schema validation;
  \item SQLite for a single-researcher prototype and PostgreSQL for concurrent deployment;
  \item a content-addressed artifact store for sources, prompts, proofs, models, and outputs;
  \item MMT for foundation-independent theory graphs \cite{rabe2016};
  \item Lean for machine-checked category theory and mathematical consequences
  \cite{leanmathlib};
  \item GATlab for computable generalized algebraic theories and theory morphisms
  \cite{gatlab};
  \item Cubical Agda or a comparable cubical system for experiments requiring computational univalence
  and higher inductive structure;
  \item SMT solvers, model finders, symbolic algebra, and numerical simulation behind typed adapters;
  \item a provider-neutral language-model interface supporting local and hosted models without making
  model-specific messages part of the domain schema.
\end{itemize}

Embeddings and semantic retrieval may propose relevant sources or alignments, but they never establish
concept identity. Accepted identity and equivalence claims require explicit profiles, maps, and
validation.

\subsection{Persistence and provenance}

The relational index stores identifiers and queryable metadata in tables for contexts, fragments,
morphisms, higher cells, pressures, meaning profiles, fragmentation sites, reflections, candidates,
validations, runs, and dependencies. Large artifacts remain immutable files addressed by a content
hash. A canonical atlas revision contains references to accepted records, not mutable copies.

Each generated claim stores:
\begin{equation}
(\mathsf{source\_set},\mathsf{skill\_version},\mathsf{model\_id},
\mathsf{prompt\_hash},\mathsf{parameters},\mathsf{seed},
\mathsf{parent\_artifacts}).
\end{equation}
This supports replay, comparison across models, and audits of whether an alleged prediction was already
present in the generation context.

\subsection{Orchestration algorithm}

\begin{lstlisting}[style=spec,language=Python,caption={Simplified orchestration loop.}]
def run_discovery(context_id: str, budget: Budget) -> RunOutcome:
    state = load_or_create_state(context_id)

    while budget.available() and not state.is_terminal():
        skill = policy.next_skill(state)
        inputs = build_typed_inputs(skill, state)

        if not skill.preconditions(inputs):
            state.block(skill, reason="precondition")
            continue

        result = execute_with_budget(skill, inputs, budget)
        validate_schema(result)
        persist_immutable(result)
        state = transition(state, result)

        if repeated_structural_failure(state):
            state.request("mutate-representation")

        if human_gate_required(state):
            return RunOutcome.awaiting_review(state)

    return synthesize_or_report_blocker(state)
\end{lstlisting}

The policy chooses the next skill by expected value of information, relevance, validation cost, and
portfolio diversity. It may run independent reflection lenses, candidate branches, and adversarial
tests concurrently. Writes to canonical state remain serialized and require schema and policy checks.

\subsection{Interfaces}

The first release should expose both a command-line interface and a small HTTP API:
\begin{lstlisting}[style=spec,caption={Minimum external interface.}]
tdr init <project>
tdr ingest <source-path> --context <id>
tdr map --context <id>
tdr reflect --site <fragmentation-id>
tdr generate --context <id> --budget <profile>
tdr validate --candidate <id>
tdr inspect --run <id>
tdr replay --run <id>
tdr export --candidate <id> --format latex|json|mmt|lean

POST /contexts
POST /contexts/{id}/sources
POST /runs
GET  /runs/{id}
GET  /sites/{id}/reflections
POST /candidates/{id}/validations
GET  /contexts/{id}/frontier
\end{lstlisting}

A research interface should visualize the local atlas, active pressures, essential fragmentation
profiles, candidate lineage, proof obligations, concrete-return paths, and Pareto frontier. It should
never imply that graph centrality or model confidence is equivalent to truth.

\subsection{Testing strategy}

Testing occurs at five levels:
\begin{enumerate}
  \item \textbf{Contract tests:} every skill accepts and emits only schema-valid records and respects
  its authority boundary.
  \item \textbf{Metamorphic tests:} renaming, equivalent presentation, unit change, and reordering do not
  materially alter accepted relevance or conclusions.
  \item \textbf{Adversarial tests:} sterile abstractions, inconsistent pushouts, false friends, prompt
  contamination, and scope erasure are rejected.
  \item \textbf{Historical reconstruction:} held-out concepts are recovered from pre-conceptual
  pressure without lexical leakage.
  \item \textbf{Prospective calibration:} expected unblocking, theorem transfer, and prediction are
  compared with realized research outcomes.
\end{enumerate}

Every benchmark contains positive and negative cases. A system that rediscovers tensor products but
also promotes every formal commonality as a deep unification has failed the intended task.

\subsection{Human review and epistemic governance}

Human approval is mandatory when:
\begin{itemize}
  \item the research context changes its goals or protected commitments;
  \item a representation mutation would migrate or discard accepted content;
  \item two lens outputs encode a genuine value or ontological choice rather than a factual question;
  \item a candidate is promoted as a substantive theory rather than a speculative research direction;
  \item an external experiment, publication claim, or irreversible action is proposed.
\end{itemize}

The interface must display uncertainty about extraction, interpretation, formalization, and validation
separately. A single confidence number would obscure where epistemic risk resides.

\subsection{Construction roadmap}

\begin{longtable}{@{}
>{\raggedright\arraybackslash}p{0.14\textwidth}
>{\raggedright\arraybackslash}p{0.24\textwidth}
>{\raggedright\arraybackslash}p{0.397\textwidth}
>{\raggedright\arraybackslash}p{0.14\textwidth}@{}}
\toprule
Phase & Build target & Measurable exit condition & Indicative duration \\
\midrule
\endfirsthead
\toprule
Phase & Build target & Measurable exit condition & Indicative duration \\
\midrule
\endhead
0 & Domain schemas, skill contract, run state, artifact provenance & Schema round-trips, replay, and authority-boundary tests pass & 3 weeks \\
1 & Atlas mapping, morphisms, source anchoring, scope tracking & Two benchmark domains mapped with at least 95\% sampled claim provenance & 5 weeks \\
2 & Pressure, meaning-profile, fragmentation, and reflection skills & Essential versus superficial fragmentation distinguished on a curated test set; every reflection emits losses and obligations & 5 weeks \\
3 & Candidate generation, relevance gates, and Pareto controller & Diverse valid portfolios produced without promoting seeded sterile completions & 4 weeks \\
4 & MMT, Lean, GATlab, solver, and simulation adapters & At least three end-to-end candidates formalized and concretely instantiated & 7 weeks \\
5 & Historical reconstruction and prospective pilot & Prespecified held-out rediscovery, invariance, unblocking, and adversarial thresholds met & 8 weeks \\
\bottomrule
\caption{Staged construction plan for a research-grade prototype.}
\label{tab:roadmap}
\end{longtable}

The first prospective pilot should use one mathematical and one physical boundary problem sharing a
structural theme. A suitable pair would be representability and obstruction discovery in mathematics,
together with state transport, coarse-graining, and spectral coercivity in gauge theory. This makes
cross-domain unification testable without assuming in advance that the same ontology governs both.

Natural-language generation is most valuable at pressure interpretation, meaning reconstruction, and
representation mutation, where the relevant abstraction may not yet have a formal name. Formal methods
become decisive once a candidate claims a particular structural role. The architecture is designed so
that neither can silently substitute for the other.

\section{Abstract-to-concrete return}

Let $U$ be an abstract candidate theory and $F_i:U\to T_i$ validated interpretations into concrete
domains. Change-of-base functors on model categories and proof translations should instantiate the
abstract construction. Where adjoints exist, Kan extensions or universal lifts may systematically
generate domain realizations; where they do not, the absence is itself a defect to be studied.

Every candidate should carry a return diagram
\begin{equation}
\begin{array}{ccc}
U & \xrightarrow{\;\varphi\;} & U' \\
\downarrow F_i && \downarrow F_i' \\
T_i & \xrightarrow{\;\varphi_i\;} & T_i',
\end{array}
\end{equation}
whose commutation states that the concrete consequence $\varphi_i$ genuinely instantiates the
abstract move $\varphi$. Approximate commutation requires an explicit error object or bound.

The return path has two purposes. First, it prevents sterile formalism. Second, it allows evidence from
the concrete domain to flow back and revise the abstract theory. The system is therefore not a one-way
pipeline from abstraction to application but a closed abductive--deductive loop.

\section{Illustration: representation systems, state connections, and mass}

Consider the proposal that a physical theory is a representation system equipped with a state
connection. Let
\begin{equation}
\pi:\mathcal{S}\longrightarrow\mathcal{R}
\end{equation}
be a fibration in which $\mathcal{R}$ is a space of descriptions and the fiber $\mathcal{S}_r$ contains
states available in representation $r$. A connection $\nabla$ transports states along descriptive
morphisms.

The atlas then generates precise pressures:
\begin{itemize}
  \item exact descriptive equivalences should induce flat or coherently path-independent transport;
  \item physical path dependence must be separated from artifacts of representation;
  \item nontrivial holonomy may encode global sectors only if it maps to known observables;
  \item coarse-graining should transform the connection naturally or with a controlled defect;
  \item any relation between curvature and interaction must be derived, not stipulated.
\end{itemize}

For Euclidean-time or transfer evolution $E_t$ with vacuum projection $P_\Omega$, a mass gap can be
expressed as a coercive estimate
\begin{equation}
\lVert E_t(1-P_\Omega)\rVert\leq C e^{-mt}.
\label{eq:gap}
\end{equation}
A relevance-guided generator may propose that $m$ measures the contraction rate of nonvacuum
relational curvature or holonomy sectors. The proposal has high potential unification value because it
connects representation transport, topology, interaction, and spectral decay. It is nevertheless only
admissible if it passes the following return tests:
\begin{enumerate}
  \item recover the usual spectral definition $H|_{\Omega^\perp}\geq m$;
  \item specify the map from curvature or holonomy data to the nonvacuum spectral subspace;
  \item establish naturality under exact changes of representation;
  \item derive a lattice transfer inequality or a continuum correlation-decay statement;
  \item produce a new bound or distinction not already contained in the spectral formulation.
\end{enumerate}

The relevance layer also detects a scope correction. Real-time quantum evolution is unitary and does
not literally dissipate nonvacuum structure. The contraction language belongs to Euclidean time,
transfer operators, heat semigroups, open systems, or coarse-grained dynamics. A modal distinction
between Lorentzian and Euclidean regimes is therefore not a technical detail but a relevance-preserving
repair: without it, the attractive unification would erase a physically essential difference.

This example shows the division of labor. Category theory identifies the possible structural bridge.
The generative function recognizes why the bridge would matter. Relevance operators prioritize the
candidate. Formal and physical validation decide whether the bridge is real.

\section{Levels of conceptual innovation}

The reasoner should label candidates by the depth at which they modify the atlas:
\begin{enumerate}
  \item \textbf{Parametric innovation:} a missing coefficient, term, boundary condition, or lemma.
  \item \textbf{Structural innovation:} a new invariant, mediator, universal object, interaction, or
  obstruction.
  \item \textbf{Representational innovation:} a new notion of state, observable, morphism, scale, or
  decomposition.
  \item \textbf{Foundational innovation:} a new criterion of identity, locality, explanation, or logic.
\end{enumerate}

High immediate unblocking relevance often occurs at Level 1. High unification and explanatory value
often occurs at Level 2. Major conceptual novelty usually requires Level 3 or 4. These dimensions must
not be conflated: a foundational proposal can be profound but irrelevant to the current bottleneck,
while a small bridge theorem can transform an active programme.

A meta-abductive trigger should activate representation mutation when:
\begin{itemize}
  \item all low-complexity completions inside the present grammar fail;
  \item successful candidates require repeated exceptions or auxiliary structures;
  \item relevance scores are unstable under equivalent presentations;
  \item concrete instantiation systematically breaks the same abstraction;
  \item independent pressure types localize at the same representational boundary.
\end{itemize}

\section{Formal properties and attainable guarantees}

The complete discovery process is not decidable or exhaustive. Several relative results are nevertheless
reasonable targets.

\begin{principle}[Transport soundness]
For an accepted theory morphism, derivability and satisfaction should be preserved in the declared
scope. A transported theorem must remain well typed and semantically valid under model reduct.
\end{principle}

\begin{principle}[Conservative completion]
For selected completion operators, sufficient conditions should guarantee that protected theorems of
the original theory are neither invalidated nor strengthened vacuously by inconsistency.
\end{principle}

\begin{principle}[Relative abductive coverage]
For a finite theory fragment, bounded complexity, and a fixed operator grammar, the system should be
able to enumerate all minimal completions up to the declared equivalence relation. This is a relative
coverage claim, not completeness over all possible concepts.
\end{principle}

\begin{principle}[Representation invariance]
If $F:D\simeq D'$ is an accepted equivalence of atlas regions, then relevance should be stable under
transport:
\begin{equation}
\mathbf{R}_Q(H;D)
\simeq
\mathbf{R}_{FQ}(FH;D'),
\end{equation}
up to explicitly declared approximation tolerances.
\end{principle}

\begin{principle}[Relevance calibration]
Relevance estimates should be empirically calibrated against downstream outcomes: proof closure,
held-out theorem transfer, successful prediction, assumption reduction, experiment design, or expert
research decisions. This is an evaluation requirement rather than a purely formal theorem.
\end{principle}

\section{Evaluation programme}

\subsection{Historical reconstruction}

The first benchmark should be counterfactual rediscovery. Construct an atlas from material that
predates or omits a known concept, then test whether the reasoner localizes the relevant pressure and
proposes the right structural role. Suitable mathematical targets include quotient objects, tensor
products, completions, localization, sheafification, spectral sequences, and representability results.
Physics targets may include state spaces, gauge connections, effective actions, order parameters,
renormalization maps, and topological sectors.

Success must be measured by the recovery of the concept's universal or operational role, not by lexical
similarity to its historical name.

\subsection{Prospective metrics}

Useful metrics include:
\begin{itemize}
  \item top-$k$ recovery of held-out universal constructions;
  \item weighted proof obligations unblocked per unit validation cost;
  \item reduction in independent assumptions or description length;
  \item number and importance of newly transported theorems;
  \item held-out predictions or conjectures produced;
  \item survival under countermodel and adversarial scope testing;
  \item stability under equivalent encodings and changes of terminology;
  \item expert judgments of explanatory and pragmatic relevance;
  \item calibration between predicted and realized research value.
\end{itemize}

Negative cases are essential. The benchmark should include elegant but sterile completions, vacuous
common abstractions, inconsistent pushouts, and high-connectivity concepts that erase decisive domain
distinctions.

\subsection{A staged implementation}

\begin{enumerate}
  \item Build a small foundation-independent theory atlas with explicit scopes, model functors, and
  provenance. MMT is a plausible knowledge substrate \cite{rabe2016}.
  \item Implement a restricted completion calculus using sketches or generalized algebraic theories.
  GATlab demonstrates that theory morphisms and symbolic models can be computational objects
  \cite{gatlab}.
  \item Connect a proof assistant for categories, functors, natural transformations, limits, adjunctions,
  and theorem transport. Lean's mathlib contains substantial category-theory infrastructure
  \cite{leanmathlib}.
  \item Add the pressure interpreter and multi-objective relevance layer.
  \item Evaluate on historical reconstruction before attempting open scientific problems.
  \item Add HoTT or cubical foundations where higher coherence is empirically necessary.
  \item Run prospective pilots at a boundary with both rich formal structure and concrete return paths.
\end{enumerate}

Category-theoretic augmentation of executable scientific models already provides evidence that
semantic model structure can support model combination and transformation \cite{halter2020}. The
proposed programme extends this direction from model manipulation to explicit opportunity detection,
relevance-guided generation, and meta-level representational revision.

\section{Limitations}

The framework has several irreducible limitations.

First, no fixed formal language can exhaust all scientifically meaningful abstractions. The atlas is a
map, not the conceptual universe. Second, relevance depends on values, aims, and opportunity costs;
these can be made explicit but not eliminated. Third, formal unification may fail to correspond to
causal or physical unity. Fourth, empirical underdetermination permits several inequivalent theories to
fit the same evidence. Fifth, formalization itself can distort a concept by privileging what is easy to
encode. Finally, generative systems inherit biases from their corpora and may repeatedly rediscover
dominant frameworks while overlooking weakly documented alternatives.

The appropriate response is not to abandon formal structure but to preserve plural representations,
track uncertainty, test invariance, retain human criticism, and treat representation mutation as a
normal part of discovery rather than as a failure of the system.

\section{Conclusion}

The proposed reasoner combines five distinct capacities:
\begin{enumerate}
  \item a higher theory atlas that represents theories, models, scopes, and transformations;
  \item a philosophical-reflective layer that identifies essential fragmentation and reconstructs
  meaning before formal completion;
  \item a categorical and homotopical calculus of structural defects and completions;
  \item a relevance-guided generative element that proposes local, bridging, mechanistic, and
  representational changes;
  \item an independent validation loop connecting abstraction to proof, computation, experiment, and
  pragmatic consequence.
\end{enumerate}

Its central innovation is to treat relevance as part of the formal discovery problem without confusing
relevance with truth. High-value concepts are not simply those that complete many diagrams. They are
those that relieve material pressures, unlock consequential work, compress explanations, unify without
erasure, transport results, generate discriminating consequences, and return to concrete practice.

The strongest defensible objective is therefore:

\begin{quote}
Given an explicit theory atlas, research context, and pressure field, identify essential conceptual
fragmentation; reconstruct the meanings and commitments at stake; generate and compare minimal
semantics-preserving transformations with high unblocking, explanatory, unifying, predictive, and
pragmatic relevance; instantiate and test their consequences; and diagnose when the highest-value move
is to revise the representation space itself.
\end{quote}

This would turn category theory from a retrospective language of unification into an operational guide
to where conceptual invention should be attempted, while the relevance-guided generative element
supplies the context-sensitive judgment that pure completion mechanics cannot provide.

\begin{thebibliography}{99}

\bibitem{hottbook}
The Univalent Foundations Program.
\newblock \emph{Homotopy Type Theory: Univalent Foundations of Mathematics}.
\newblock Institute for Advanced Study, 2013.
\newblock \url{https://homotopytypetheory.org/book/}

\bibitem{rabe2016}
Florian Rabe.
\newblock MMT: A foundation-independent approach to formal knowledge.
\newblock 2016.
\newblock \url{https://uniformal.github.io/doc/philosophy/articles/mmt.pdf}

\bibitem{goguenburstall1992}
Joseph A. Goguen and Rod M. Burstall.
\newblock Institutions: Abstract model theory for specification and programming.
\newblock \emph{Journal of the ACM}, 39(1):95--146, 1992.

\bibitem{fongspivak2019}
Brendan Fong and David I. Spivak.
\newblock \emph{An Invitation to Applied Category Theory: Seven Sketches in Compositionality}.
\newblock Cambridge University Press, 2019.
\newblock \url{https://arxiv.org/abs/1803.05316}

\bibitem{gatlab}
Owen Lynch, James Fairbanks, Evan Patterson, and collaborators.
\newblock GATlab: Modeling and programming with generalized algebraic theories.
\newblock 2024.
\newblock \url{https://arxiv.org/abs/2404.04837}

\bibitem{halter2020}
Micah Halter, Christine Herlihy, and James Fairbanks.
\newblock A compositional framework for scientific model augmentation.
\newblock \emph{Electronic Proceedings in Theoretical Computer Science}, 323:172--182, 2020.
\newblock \url{https://arxiv.org/abs/1907.03536}

\bibitem{leanmathlib}
Lean Community.
\newblock Category theory in mathlib.
\newblock \url{https://leanprover-community.github.io/theories/category_theory.html}

\bibitem{kitcher1981}
Philip Kitcher.
\newblock Explanatory unification.
\newblock \emph{Philosophy of Science}, 48(4):507--531, 1981.

\bibitem{harman1965}
Gilbert H. Harman.
\newblock The inference to the best explanation.
\newblock \emph{The Philosophical Review}, 74(1):88--95, 1965.

\bibitem{rissanen1978}
Jorma Rissanen.
\newblock Modeling by shortest data description.
\newblock \emph{Automatica}, 14(5):465--471, 1978.

\bibitem{kuhn1962}
Thomas S. Kuhn.
\newblock \emph{The Structure of Scientific Revolutions}.
\newblock University of Chicago Press, 1962.

\bibitem{lakatos1976}
Imre Lakatos.
\newblock \emph{Proofs and Refutations: The Logic of Mathematical Discovery}.
\newblock Cambridge University Press, 1976.

\bibitem{brandom1994}
Robert B. Brandom.
\newblock \emph{Making It Explicit: Reasoning, Representing, and Discursive Commitment}.
\newblock Harvard University Press, 1994.

\bibitem{cassirer1910}
Ernst Cassirer.
\newblock \emph{Substance and Function}.
\newblock 1910; English translation, Open Court, 1923.

\end{thebibliography}

\end{document}
